
\subsubsection{2.1 N-gram Contamination Measurement}

Systematic study of training data contamination began with GPT-3 \citep{Brownetal2020}, which reported n-gram overlap between test sets and training data and applied decontamination filters. That analysis was limited to a closed training corpus and did not provide per-sub-task breakdowns.

The most directly relevant prior work is WIMBD (What's In My Big Data?) \citep{Elazaretal2023}, which provides open-source tooling for 13-gram containment analysis on The Pile v1. WIMBD demonstrates that contamination rates vary across MMLU sub-tasks within The Pile and establishes 13-gram containment as a standard measurement approach. The present work adopts the same 13-gram methodology, validates the analysis pipeline against WIMBD published rates (Spearman ρ=0.721), and extends the analysis from a single corpus to a three-corpus comparison. WIMBD covers The Pile only; C4 and RedPajama remain uncharacterized by WIMBD.

The GPT-4 Technical Report [OpenAI, 2023] reports contamination analysis using Jaccard similarity for select benchmarks against a closed training corpus. Riddell et al. [2024] apply n-gram contamination analysis to code generation benchmarks, demonstrating the generality of the methodology. Burnell et al. [2023] call for contamination-aware evaluation and note the absence of systematic cross-corpus contamination studies.

\subsubsection{2.2 Corpus Documentation and Curation}

Gao et al. [2020] introduce The Pile v1, a diverse 825 GB corpus assembled from 22 sub-sources including PubMed Central, ArXiv, FreeLaw, and Common Crawl. The deliberate inclusion of domain-aligned academic sources is the property expected to produce the highest benchmark contamination rates among the three corpora studied. Dodge et al. [2021] document C4 (Colossal Clean Crawled Corpus), derived from Common Crawl with aggressive quality filtering. TogetherComputer [2023] introduce RedPajama-v1 as an open reproduction of LLaMA's training data, combining Common Crawl, C4, GitHub, Wikipedia, Books, ArXiv, and StackExchange.

\subsubsection{2.3 Contamination and Performance Inflation}

Magar and Schwartz [2022] provide a theoretical framework connecting training data contamination to performance inflation. Their work motivates contamination measurement but does not provide an empirical cross-corpus contamination study. The present work provides that measurement.

\subsubsection{2.4 Positioning}

This work extends WIMBD's single-corpus analysis to three corpora; characterizes C4 and RedPajama contamination profiles that corpus documentation papers leave unmeasured; and provides sub-task-level contamination estimates enabling corpus-vs-corpus comparison. Unlike WIMBD, which characterizes one corpus against individual datasets in isolation, this work produces a unified matrix enabling systematic, directly comparable contamination estimates across corpora at sub-task granularity.

